% Faithful formal-publication transcription: Section3, pp23:15--23:24. Historical intuition is not needed as substitute for the full formal proofs below.
\section*{3 Randomness-efficient low-complexity hash functions}
In this section, we present constructions of randomness-efficient low-complexity hash functions with various parameters and properties.
In Section3.1, we show almost universal hash functions can be constructed from 1-detectors (Lemma3.2). In Section3.2 we give a derandomized version of construction of 1-detectors from [13] based on high-girth graphs (Lemma3.3), and use that to construct a low-complexity hash function with poly-logarithmic seed-length and $n/\operatorname{polylog}(n)$ output length (Theorem3.4).
In Section3.3, we show how to reduce the output length from $n/\operatorname{polylog}(n)$ to $\operatorname{polylog}(n)$, by composing almost universal hash families. Finally, in Section3.4 and3.5, we establish the uniformity and explicitness of our constructions, which are essential for our applications to hardness magnification and PRF constructions.
\paragraph{Notation.} In this section, for a construction $H$ with parameters $a,b,c$, we will use $H_{\hat a,\hat b,\hat c}$ to denote the specific construction $H$ with the parameters specified to $\hat a,\hat b$ and $\hat c$. When the parameters are obvious from the context, we often omit the subscripts and simply write it as $H$.
\subsection*{3.1 General construction from 1-detectors}
We first give a general construction of almost universal hash functions from 1-detectors, using the following sampling procedure.
\begin{authorclaim}{Lemma3.1}
For all integers $n,b$ and $r$ such that $b\le n$ and $\max\{10n/b,10\log\log n\}\le r\le n$, there is a distribution $D^{\mathrm{samp}}_{n,b,r}$ supported on $[n]^b$ samplable by $O(r\log(n/b))$ bits, such that the following conditions hold:
\begin{enumerate}
\item (Hitting Condition). For all $x\in\mathbb F_2^n$ with $|x|\ge r$, it holds that
\begin{equation}\tag{1}
\Pr_{(w_0,\ldots,w_{b-1})\leftarrow D^{\mathrm{samp}}_{n,b,r}}\left[\bigwedge_{j\in[b]}[x_{w_j}=0]\right]\le2\exp\left[-\frac{br}{2n}\right].
\end{equation}
\item (Ordering) For every $(w_0,\ldots,w_{b-1})\in\operatorname{supp}(D^{\mathrm{samp}}_{n,b,r})$, it holds that $w_0<w_1<\cdots<w_{b-1}$.
\item (Explicitness) There are algorithms $A^{\mathrm{samp}}_{n,b,r}(v,i)$ and $B^{\mathrm{samp}}_{n,b,r}(v,j)$ running in $\operatorname{poly}(\log n,r)$ time such that
\[A^{\mathrm{samp}}_{n,b,r}(v,i)\text{ outputs }\begin{cases}j&\text{if }w_j=i\text{ for some }j\in[b],\\\bot&\text{otherwise;}\end{cases}\]
\[B^{\mathrm{samp}}_{n,b,r}(v,j)\text{ outputs }w_j,\]
where $(w_0,\ldots,w_{b-1})\in\operatorname{supp}(D^{\mathrm{samp}}_{n,b,r})$ is the vector corresponding to the seed $v$.
\end{enumerate}
\end{authorclaim}
\begin{proof}
We firstly describe a sampling procedure that satisfies Equation(1) but has a long seed length and then reduce the seed length to $O(r\log(n/b))$ bits using the explicit $k$-wise $\varepsilon$-dependent distribution from Theorem2.6. The sampling procedure is simple: we first partition $[n]$ into $b$ consecutive groups $g_0,g_1,\ldots,g_{b-1}$ of size either $\lfloor n/b\rfloor$ or $\lceil n/b\rceil$; we then uniformly choose an $i_j\in g_j$ for each $j\in[b]$; finally we output the tuple $(i_0,i_1,\ldots,i_{b-1})$.
For any $x\in\mathbb F_2^n$ with Hamming weight at least $r$, for some $\ell\le r$, there are $\ell$ groups such that there are at least $r$ 1-indices (i.e., the corresponding bits of $x$ are 1) in these groups. Assume that there are $a_j$ 1-indices in the $j$th among these $\ell$ groups. The probability that none of the 1-indices is sampled is at most
\begin{equation}\tag{2}
\begin{aligned}
\prod_{j=1}^\ell\left[1-\frac{a_jb}{2n}\right]
&\le\exp\left(\sum_{j=1}^\ell\log\left[1-\frac{a_jb}{2n}\right]\right)\\
&\le\exp\left(-\sum_{j=1}^\ell\frac{a_jb}{2n}\right)\le\exp\left[-\frac{br}{2n}\right].
\end{aligned}
\end{equation}
A trivial implementation of the sampling procedure above needs seed length $O(b\log(n/b))$: for each group, we sample a random binary string of length $t\triangleq\lceil\log(n/b)\rceil$ indicating the index we want to choose. Therefore, we would need to sample $bt$ bits in total. The observation here is that the argument in (2) only involves $\ell\le r$ groups. Hence, the analysis still works if we sample those required $bt$ bits from an $rt$-wise $\exp(-br/(2n))$-dependent distribution with alphabet 2. Note that this incurs an additional error of $\exp(-br/(2n))$.
By Theorem2.6, the distribution is samplable with seed length
\[O(rt+\log\log(bt)+br/(2n))=O(r\log(n/b)),\]
and the algorithms $A^{\mathrm{samp}}_{n,b,r}$ and $B^{\mathrm{samp}}_{n,b,r}$ can be constructed straightforwardly from the algorithm $B$ in Theorem2.6.\footnote{To compute $A^{\mathrm{samp}}_{n,b,r}(v,i)$ for some $i$ belonging to the group $g_a$, we only need to check if $i$ is the selected element of the group $g_a$, and then output $a$ if the answer is yes, and output $\bot$ otherwise.}
\end{proof}
Now we are ready to give the construction of our hash function from 1-detectors. Recall that their definitions are given in Section2.3.
\begin{authorclaim}{Construction1 (Hash function $H^L$ from 1-detectors $L$)}
Let $n$ be the input length and $b(n)=o(n)$ be a parameter. Given a randomized 1-detector $L=\{L_n\}_{n\ge1}$ with range $\hat r$, we construct a hash function $H_b^L=\{H_n\}_{n\ge1}$ as follows.
\begin{itemize}
\item Let $D^{\mathrm{samp}}$ be the distribution in Lemma3.1 with parameters $n,b$ and $r=\hat r$. For each $L$ in the support of $L$ and each $D=(j_0,\ldots,j_{b-1})$ in the support of $D^{\mathrm{samp}}$, we define a function $h_{L,D}$ as
\[h_{L,D}(x)\triangleq L(x)\Vert x_{j_0}\Vert x_{j_1}\Vert\cdots\Vert x_{j_{b-1}}.\]
\item $H_n$ is then defined to be the distribution generated as follows: sample $L\leftarrow L_n$ and $D\leftarrow D^{\mathrm{samp}}$, and then output $h_{L,D}$.
\end{itemize}
\end{authorclaim}
We show that Construction1 indeed gives almost universal hash functions.
\begin{authorclaim}{Lemma3.2 ($H^L$ is a linear almost universal hash)}
Let $L=\{L_n\}_{n\ge1}$ be an $\hat m$-output $\hat s$-samplable randomized 1-detector with error $\hat\varepsilon$ and range $\hat r$, such that $\hat m=o(n)$ and $\hat r=\omega(\log n)$. Let $b=b(n)=o(n)$ be a parameter and $r=r(n)=\hat r(n)$. The followings hold for $H_b^L$:
\begin{enumerate}
\item $H_b^L$ has output length $m=\hat m+b$.
\item $H_b^L$ is $(\hat s+\hat r\log(n/b))$-samplable.
\item $H_b^L$ is a linear $\varepsilon$-almost universal hash function, where $\varepsilon\triangleq\hat\varepsilon+\exp(-\Omega(br/n))$.
\end{enumerate}
\end{authorclaim}
\begin{proof}
The first two items follow directly from the definition of $H^L$ from Construction1. So we will only establish the last item. We need to show that for any two inputs $x_1\ne x_2$ from $\{0,1\}^n$, their hash values collide with probability at most $\varepsilon$.
Consider the following two cases depending on whether the Hamming distance between $x_1$ and $x_2$ is small or large:
\begin{itemize}
\item If $0<\Delta(x_1,x_2)\le r$, then by the definition of randomized 1-detector (Definition2.2), $L(x_1)$ equals to $L(x_2)$ with probability at most $\hat\varepsilon$ for $L\leftarrow L_n$ (note that $L$ is a linear function over $\mathbb F_2$). Since the hash values contain the output of the 1-detector as the first $\hat m$ bits, the hash values collide with probability at most $\hat\varepsilon$.
\item If $\Delta(x_1,x_2)>r$, then $y\triangleq x_1\oplus x_2$ has Hamming weight at least $r$. By Lemma3.1, a random $D\leftarrow D^{\mathrm{samp}}_{n,b,r}$ fails to contain all 1-indices of $y$ with probability at most $\exp(-\Omega(br/n))$, which means that the last $b$ bits of the hash values of $x_1$ and $x_2$ collide with probability at most $\exp(-\Omega(br/n))$.
\end{itemize}
It follows immediately that the collision probability is at most $\varepsilon=\hat\varepsilon+\exp(-\Omega(br/n))$.
\end{proof}
\subsection*{3.2 Randomness-efficient low-complexity 1-detectors}
Now we will present the construction of randomness-efficient low-complexity 1-detectors. Applying Lemma3.2, this construction yields a hash function with inverse-super-polynomial collision probability, short seed length, and moderately large shrinkage.
\begin{authorclaim}{Construction2 (Construction of the 1-detector $L^{\mathrm{hg}}$ from high-girth graphs)}
Let $k=k(n)=n^{o(1)}$ be the shrinkage parameter, $\log n\le r(n)\le n/2$ be the range. Let $\gamma=\gamma(n)\triangleq\log^2n/\log k$. $L^{\mathrm{hg}}_{k,r}=\{L_n\}_{n\ge1}$ is constructed as follows.
\begin{itemize}
\item Let $n'$ be the smallest power of two that is larger than $n$. Let $G_{m,n'}$ be the regular graph with $m=\Theta(n'/k)$ vertices, $n'$ edges, and girth $g=\Theta(\log n/\log k)$ from Theorem2.12. Let $D=\{D_{n'}\}$ be the explicit family of $r$-wise $2^{-\gamma}$-dependent permutation in Theorem2.8. Let $\Gamma(i)$ be the set of indices of edges incident to the $i$th vertex in $G_{m,n'}$. For each permutation $\sigma\in\operatorname{supp}(D_{n'})$, we define a function $L_\sigma:\{0,1\}^n\to\{0,1\}^m$ as
\[L_\sigma(x)_i\triangleq\bigoplus_{j\in[n]\text{ s.t. }\sigma(j)\in\Gamma(i)}x_j\qquad\forall i\in[m].\]
\item $L_n$ is then defined to be the distribution generated as follows: sample $\sigma\leftarrow D_{n'}$ and then output $L_\sigma$.
\end{itemize}
\end{authorclaim}
\begin{authorclaim}{Lemma3.3 ($L^{\mathrm{hg}}$ is a 1-detector)}
Let $k=k(n)=n^{o(1)}$ be the shrinkage parameter, $\log n\le r(n)\le n/2$ be the range. The followings hold for $L^{\mathrm{hg}}_{k,r}=\{L_n\}_{n\ge1}$:
\begin{enumerate}
\item $L^{\mathrm{hg}}_{k,r}$ has output size $\Theta(n/k)$ and seed length $O(r\log n)$.
\item Every $L\in L_n$ can be computed by depth-1 $\mathsf{CC}^0[2]$ circuits of wire complexity at most $2n$, or $B_2$ circuits of size $2n$ and depth $\log k+O(1)$.
\item $L^{\mathrm{hg}}_{k,r}$ is a randomized 1-detector with range $r$ and error $\exp(-\Omega(\log^2n/\log k))$.
\end{enumerate}
\end{authorclaim}
\begin{proof}
Recall that $\gamma=\log^2n/\log k$ from Construction2. For the first item, the output size follows directly from the definition of $L^{\mathrm{hg}}_{k,r}$. The seed length of $L^{\mathrm{hg}}_{k,r}$ equals the seed length of the required $r$-wise $2^{-\gamma}$-dependent permutation, which is $O(r\log n+\gamma)=O(r\log n)$ from Theorem2.8 (note that $r\ge\log n$). The second item follows directly from the definition of $L_\sigma$ in Construction2 together with the fact that $G_{m,n'}$ has maximum degree $O(k)$ (since it is regular).
In the rest of the proof we establish the third item. We need to show that for any non-zero $x\in\mathbb F_2^n$, the probability that $L(x)=0$ with $L\leftarrow L_n$ is at most $\exp(-\Omega(\gamma))$. For any input $x$ with Hamming weight $0<|x|<g$ (where $g$ is the girth of $G_{m,n}$), we must have $L_\sigma(x)\ne0$ for all $\sigma\in D_{n'}$, since otherwise we can extract from the edges $\{\sigma(i)\mid x_i=1\wedge i\in[n]\}$ a cycle of size length less than $g$. Now we consider the case when $g\le|x|\le r$. Let $i_1,i_2,\ldots,i_{|x|}$ be the indices corresponding to the 1's in $x$ (i.e., $x_{i_j}=1$ for all $j$). Since our construction is linear, only these bits can influence the output. Also, by the $r$-wise $2^{-\gamma}$-dependence of $D_{n'}$, we have
\begin{align*}
\Pr_{\sigma\leftarrow D_{n'}}[L_\sigma(x)=0]
&=\Pr_{\sigma\leftarrow D_{n'}}\left[\bigwedge_{i'\in[m]}[|\{i_j\mid\sigma(i_j)\in\Gamma(i')\}|\text{ is even}]\right]\\
&\le\Pr_{S\subseteq[n'],\,|S|=|x|}\left[\bigwedge_{i'\in[m]}[|S\cap\Gamma(i')|\text{ is even}]\right]+2^{-\gamma}\\
&=\mathbb E_{\substack{S\subseteq[n']\\|S|=|x|-\alpha}}\left[\Pr_{\substack{S'\subseteq[n']\setminus S\\|S'|=\alpha}}\left[\bigwedge_{i'\in[m]}[|(S\cup S')\cap\Gamma(i')|\text{ is even}]\right]\right]+2^{-\gamma},
\end{align*}
where $\alpha\triangleq\lceil g/3\rceil$.
The final observation is that for any fixed $S$, there should be at most one $S'$ satisfying the condition in the summation. Indeed, if two sets $S'_1$ and $S'_2$ satisfy the condition at the same time, then the symmetric difference of them contains a cycle in $G_{m,n'}$ of length $2\alpha<g$ for sufficiently large $n$, contracting the fact that $G_{m,n'}$ has girth $g$. In particular, it means for every fixed $S$, we have
\[\Pr_{\substack{S'\subseteq[n']\setminus S\\|S'|=\alpha}}\left[\bigwedge_{i'\in[m]}[|(S\cup S')\cap\Gamma(i')|\text{ is even}]\right]\le\frac1{\binom{n'-|x|+\alpha}{\alpha}}.\]
Since $\alpha=\lceil g/3\rceil=\Theta(\log n/\log k)$, putting everything together, for sufficiently large $n$ we have
\[\Pr_{\sigma\leftarrow D_n}[L_\sigma(x)=0]\le\frac1{\binom{n'-|x|+\alpha}{\alpha}}+2^{-\gamma}=\exp(-\Omega(\gamma)).\]
\end{proof}
Plugging $L^{\mathrm{hg}}$ into Construction1 with appropriately chosen parameters, we immediately obtain a linear hash $H^{\mathrm{md}}$ with poly-logarithmic seed length and shrinkage as follows. (Here md stands for moderate, meaning that the hash has moderate (poly-logarithmic) shrinkage.)
\begin{authorclaim}{Construction3 (Construction of hash $H^{\mathrm{md}}$ with poly-logarithmic shrinkage)}
Let $\beta\in(0,2]$ be a constant and $\omega(\log n)\le\delta(n)\le O(\log^2n/\log\log n)$ be the error parameter. $H^{\mathrm{md}}_{\beta,\delta}=\{H_n\}_{n\ge1}$ is constructed as follows.
\begin{itemize}
\item Setting $k=\log^\beta n$ and $r=\delta\log^\beta n$, $L^{\mathrm{hg}}=L^{\mathrm{hg}}_{k,r}$ (from Construction2) is an $O(\delta\log^{\beta+1}n)$-samplable $\Theta(n/\log^\beta n)$-output randomized 1-detector with range $r$ and error $\exp(-\Omega(\delta))$.
\item Setting $b=n/\log^\beta n$, $H^{\mathrm{md}}$ is now defined to be $H_b^{L^{\mathrm{hg}}}$ (from Construction1).
\end{itemize}
\end{authorclaim}
The following theorem follows directly from the Construction3, Lemma3.2, and Lemma3.3.
\begin{authorclaim}{Theorem3.4 (Properties of the intermediate hash construction $H^{\mathrm{md}}$)}
Let $\beta\in(0,2]$ be the shrinkage parameter and $\omega(\log n)\le\delta(n)\le O(\log^2n/\log\log n)$ be the error parameter. The followings hold for $H^{\mathrm{md}}_{\beta,\delta}=\{H_n\}_{n\ge1}$:
\begin{enumerate}
\item $H^{\mathrm{md}}_{\beta,\delta}$ is a linear $O(\delta\log^{\beta+1}n)$-samplable $\Theta(n/\log^\beta n)$-output $\exp(-\Omega(\delta))$-almost universal hash function.
\item Every $H\in H_n$ can be computed either by a depth-1 $\mathsf{CC}^0[2]$ circuit of wire complexity $2n$, or by a $B_2$ circuit of size $2n$ and depth $\beta\log\log n+O(1)$.
\end{enumerate}
\end{authorclaim}

\subsection*{3.3 Shrinkage reduction of hash function}
Now we reduce the output length of the hash function $H^{\mathrm{md}}$ in Construction3 from $n/\operatorname{polylog}(n)$ to $\operatorname{polylog}(n)$ with little overhead in its circuit complexity and seed length, which is crucial for our applications. The idea is simple: we composite it with a hash with large shrinkage, short seed length, and relatively larger circuit complexity. In particular, we can use the following hash construction $H^{\mathrm{expw}}$ of Chen, Jin, and Williams [9]. (Here expw stands for expander walk.)
\begin{authorclaim}{Construction4 (Hash function $H^{\mathrm{expw}}$ from expander walk [9])}
Let $\ell=\ell(n)$ be the output length. $H^{\mathrm{expw}}_\ell=\{H_n\}_{\ge n}$ is constructed as follows.
\begin{itemize}
\item Let $\{S_n\}_{n\ge1}$ be a family of $0.1$-biased set from Theorem2.4 and $\{G_n\}_{n\ge1}$ be the family of strongly explicit expanders from Theorem2.10. Assume that $S_n=\{w_0,w_1,\ldots,w_{t-1}\}$ for $t=\widetilde O(n^2)$.
\item For each walk $v=(v_0,v_1,\ldots,v_{\ell-1})\in[t]^\ell$ of length $\ell$ on $G_t$, we define a hash function $h_v:\{0,1\}^n\to\{0,1\}^\ell$ as
\[h_v(x)\triangleq\langle w_{v_0},x\rangle\Vert\langle w_{v_1},x\rangle\Vert\cdots\Vert\langle w_{v_{\ell-1}},x\rangle,\]
where $\langle a,b\rangle$ denotes the inner product over $\mathbb F_2$.
\item $H_n$ is then defined to be the distribution generated as follows: samples a random walk $v$ of length $\ell$ on $G_t$ uniformly at random, then outputs $h_v$.
\end{itemize}
\end{authorclaim}
\begin{authorclaim}{Lemma3.5 (Properties of the hash construction $H^{\mathrm{expw}}$ from [9])}
Let $\ell=\ell(n)$ be the output length. The followings hold for $H^{\mathrm{expw}}_\ell=\{H_n\}_{n\ge1}$:
\begin{enumerate}
\item Every $H\in\operatorname{supp}(H_n)$ can implemented by a depth-1 $\mathsf{CC}^0[2]$ circuit of wire complexity $n\ell$ or a $B_2$ circuit of size $n\ell$ and depth $\log n+O(1)$.
\item $H^{\mathrm{expw}}_\ell$ is a linear $O(\log n+\ell)$-samplable $\ell$-output $\exp(-\Omega(\ell))$-almost universal hash function.
\end{enumerate}
\end{authorclaim}
\begin{proof}
The first item follows directly from the definition of $h_v$ from Construction4. In the following we establish the second item.
The linearity, seed length, and output length are straightforward to verify, thus we only analyze the collision probability. Since the hash function is linear, we only need to show that for any non-zero $x\in\mathbb F_2^n$, $\Pr_{h\leftarrow H_n}[h(x)=0]\le\exp(-\Omega(\ell))$. By Theorem2.4, we know that for any non-zero $x$, at least $0.4$ fraction of vectors $w\in S_n$ satisfies $\langle w,x\rangle=1$. Then by Theorem2.11, a random walk of length $\ell$ will find one of such $w$ with probability $1-\exp(-\Omega(\ell))$. This immediately implies that $h(x)\ne0$ with probability at most $\exp(-\Omega(\ell))$ for $h\leftarrow H_n$.
\end{proof}
Next we formally define the composition of two hash function families.
\begin{authorclaim}{Definition3.6 (Composition of hash families)}
Let $F=\{F_n\}_{n\ge1}$ and $G=\{G_n\}_{n\ge1}$ be two families of hash functions. The composition of $F$ and $G$ is defined as $F\circ G=\{(F\circ G)_n\}_{n\ge1}$, where $(F\circ G)_n$ is the following distribution: let $m=m(n)$ be the output length of $G$, we sample $f\leftarrow F_m$ and $g\leftarrow G_n$, and then output $f\circ g$.
\end{authorclaim}
The following proposition is crucial for the analysis of our final hash construction.
\begin{authorclaim}{Proposition3.7}
The composition $H'\circ H$ of an $\varepsilon_1$-almost universal hash function $H'$ and an $\varepsilon_2$-almost universal hash function $H$ is an $(\varepsilon_2(n)+\varepsilon_1(m(n)))$-almost universal hash function, where $m(n)$ is the output length of $H$. Moreover, $H'\circ H$ is linear if both of $H'$ and $H$ are linear.
\end{authorclaim}
The following is a simple corollary of Proposition3.7.
\begin{authorclaim}{Corollary3.8}
For any $t\ge2$, non-increasing $\varepsilon=\varepsilon(n)$, and output length parameter $\ell(n)\le n$, the $t$th order composition of an $\ell$-output $\varepsilon$-almost universal hash $H$ with itself, denoted by $H^{\circ t}$, is also a hash function with collision probability $t\cdot\varepsilon(\hat\ell)$, where $\hat\ell=\ell\circ\ell\circ\cdots\circ\ell(n)$ ($t-1$ times) is the input of the outer-most hash. Moreover, $H^{\circ t}$ is linear if $H$ is linear.
\end{authorclaim}
We are now ready to specify our final hash construction $H^{\mathrm{final}}$ with $\operatorname{polylog}(n)$ output length.
\begin{authorclaim}{Construction5 (Construction of $H^{\mathrm{final}}$ with $\operatorname{polylog}(n)$ output length)}
Let $\beta\in(0,2]$ be a parameter and $\omega(\log n)\le\ell(n)\le O(\log^2n/\log\log n)$ be a non-decreasing function. We define
\[H^{\mathrm{final}}_{\beta,\ell}\triangleq H^{\mathrm{expw}}_\ell\circ(H^{\mathrm{md}}_{\beta,\ell})^{\circ\lceil2/\beta\rceil}.\]
\end{authorclaim}
We analyze the properties of $H^{\mathrm{final}}$ constructed above by combining the composition proposition (Proposition3.7 and Corollary3.8) together with properties of $H^{\mathrm{expw}}$ (Lemma3.5) and $H^{\mathrm{md}}$ (Theorem3.4).
\begin{authorclaim}{Theorem3.9 (Properties of the final hash construction $H^{\mathrm{final}}$)}
Let $\beta\in(0,2]$ be a constant and $\omega(\log n)\le\ell(n)\le O(\log^2n/\log\log n)$ be a non-decreasing function. The followings hold for $H^{\mathrm{final}}_{\beta,\ell}=\{H_n\}_{n\ge1}$:
\begin{enumerate}
\item $H^{\mathrm{final}}_{\beta,\ell}$ is a linear $O(\ell\log^{\beta+1}n)$-samplable $\Theta(\hat\ell)$-output $\exp(-\Omega(\hat\ell))$-almost universal hash function, where $\hat\ell=\ell(\Theta(n/\log^{\beta\cdot\lceil2/\beta\rceil}n))$.
\item Every $H\in H_n$ can be computed by depth-$(1+\lceil2/\beta\rceil)$ $\mathsf{CC}^0[2]$ circuits of wire complexity $2n+O(n\hat\ell/\log^2n)$, or by $B_2$ circuits of size $2n+O(n\hat\ell/\log^2n)$ and depth $\log n+O(1)$.
\end{enumerate}
\end{authorclaim}
\begin{proof}
Let $H^{\mathrm{md}}=H^{\mathrm{md}}_{\beta,\ell}$ and $H^{\mathrm{expw}}=H^{\mathrm{expw}}_\ell$ be the hash functions described in Construction5. We also use $H_n^{\mathrm{md}}$ and $H_n^{\mathrm{expw}}$ to denote the distributions over $n$-input hash functions in $H^{\mathrm{md}}$ and $H^{\mathrm{expw}}$, respectively.
\paragraph{$H^{\mathrm{final}}$ is almost universal with $\operatorname{polylog}(n)$ output length.}
Let $\hat m$ be the output length of $(H^{\mathrm{md}})^{\circ\lceil2/\beta\rceil}$, and $\hat\ell\triangleq\ell(\hat m)=\ell(\Theta(n/\log^{\beta\cdot\lceil2/\beta\rceil}))$. By a simple induction and Theorem3.4 we know that for any $1\le d\le\lceil2/\beta\rceil$, $(H^{\mathrm{md}})^{\circ d}$ is a linear $O(\ell\log^{\beta+1}n)$-samplable $\Theta(n/\log^{\beta d}n)$-output $\exp(-\Omega(\hat\ell))$-almost universal hash function. In particular, when $d=\lceil2/\beta\rceil$, we know that $(H^{\mathrm{md}})^{\circ\lceil2/\beta\rceil}$ is a linear $O(\ell\log^{\beta+1}n)$-samplable $\Theta(n/\log^{\beta\cdot\lceil2/\beta\rceil}n)$-output $\exp(-\Omega(\hat\ell))$-almost universal hash function.
Now apply Proposition3.7 and Lemma3.5, we know that $H^{\mathrm{final}}=H^{\mathrm{expw}}\circ(H^{\mathrm{md}})^{\circ\lceil2/\beta\rceil}$ is a linear $O(\ell\log^{\beta+1}n)$-samplable $\Theta(\hat\ell)$-output $\exp(-\Omega(\hat\ell))$-almost universal hash function.
\paragraph{Complexity with $\mathsf{CC}^0[2]$ circuits.}
By Theorem3.4, every $H\in\operatorname{supp}(H_n^{\mathrm{md}})$ can be computed by a depth-1 $\mathsf{CC}^0[2]$ circuit with wire complexity $2n$. Also, by Lemma3.5, every $H\in\operatorname{supp}(H_n^{\mathrm{expw}})$ can be computed by a depth-1 $\mathsf{CC}^0[2]$ circuit with wire complexity $n\ell$. Since the output length of $(H^{\mathrm{md}})^{\circ d}$ is $\Theta(n/\log^{\beta d})$, the total wire complexity of a hash function from the support of $H^{\mathrm{expw}}\circ(H^{\mathrm{md}})^{\circ\lceil2/\beta\rceil}$ can be bounded by
\[2n+O\left(\sum_{d=1}^{\lceil2/\beta\rceil-1}\frac n{\log^{\beta d}n}\right)+\frac n{\log^{\beta\cdot\lceil2/\beta\rceil}n}\cdot\hat\ell=2n+O\left(\frac{n\hat\ell}{\log^2n}\right).\]
The depth complexity can be verified straightforwardly.
\paragraph{Complexity with $B_2$ circuits.}
We only analyze the depth since the analysis of the size complexity is similar to the analysis above. By Theorem3.4 and Lemma3.5, every $H\in\operatorname{supp}(H_n^{\mathrm{md}})$ can be computed by a $B_2$ circuit with depth $\beta\log\log n+O(1)$, and every $H\in\operatorname{supp}(H_n^{\mathrm{expw}})$ can be computed by a $B_2$ circuit with depth $\log n+O(1)$. Therefore the total depth of a hash function from the support of $H^{\mathrm{expw}}\circ(H^{\mathrm{md}})^{\circ\lceil2/\beta\rceil}$ is at most
\[\sum_{d=0}^{\lceil2/\beta\rceil-1}\left[\beta\log\log\left(\Theta\left(\frac n{\log^{\beta d}n}\right)\right)+O(1)\right]+\log\left(\Theta\left(\frac n{\log^{\beta\cdot\lceil2/\beta\rceil}n}\right)\right)=\log n+O(1).\]
\end{proof}
\subsection*{3.4 Explicitness of our construction}
Apart from the seed length, shrinkage and collision probability, the explicitness of the a hash function is also critical for many applications. That is, given a seed $v$ (which is usually much shorter than the input, say $|v|=\operatorname{polylog}(n)$), whether we can obtain information about the hash function corresponding to the seed $v$ efficiently. In this section we show that our hash constructions are indeed explicit in a strong sense, which is crucial for our application to hardness magnification.
\begin{authorclaim}{Definition3.10 (Locally explicit hash function)}
A family of $\operatorname{polylog}(n)$-samplable distributions over $m$-output linear functions (e.g., linear hash functions or 1-detectors) $H=\{H_n\}_{n\ge1}$ is said to be locally explicit if each input bit only influences $\operatorname{polylog}(n)$ output bits, and there exists an algorithm $A(n,v,i)$ running in $\operatorname{polylog}(n)$ time that returns the list of output bits influenced by the $i$th input bit in the hash function corresponding to the seed $v$.
\end{authorclaim}
Note that any linear function can be realized as $x\mapsto Mx$ for some transformation matrix $M$ over $\mathbb F_2$. Clearly, a locally explicit linear hash has a sparse transformation matrix, and we can efficiently list all 1-entries in any column of it. By performing a sparse matrix multiplication, we can immediately obtain the following proposition.
\begin{authorclaim}{Proposition3.11}
The composition of two locally explicit hash families is locally explicit.
\end{authorclaim}
Now we verify that our constructions of the randomness-efficient low-complexity hash functions are indeed locally explicit.
\begin{authorclaim}{Theorem3.12}
Let $k=k(n)=n^{o(1)}$ be the shrinkage parameter and $\log n\le r(n)\le n/2$ be the range. The 1-detector $L^{\mathrm{hg}}_{k,r}$ in Construction2 is locally explicit.
\end{authorclaim}
\begin{proof}
Recall that the 1-detector in Construction2 consists of two components: a depth-1 $\mathsf{CC}^0[2]$ circuit whose topology is determined by the high-girth graph from Theorem2.12, and a random permutation of the input bits according to the $k$-wise almost independent permutation from Theorem2.8. Given the seed $v$ of the 1-detector (i.e., the seed for the permutation) and an index $i$. According to Theorem2.8, we can obtain $\sigma_v(i)$ in $\operatorname{poly}(\log n,|v|)$ time, where $\sigma_v$ is the permutation corresponding to the seed $v$. To find all the output bits influenced by the $i$th input bit, we only need to find the two endpoints of the $\sigma_v(i)$-th edge in the high-girth graph, which can be done in $\operatorname{polylog}(n)$ time by Theorem2.12.
\end{proof}
\begin{authorclaim}{Theorem3.13}
Let $\beta>0$ be the shrinkage parameter and $\delta=\delta(n)\le O(\log^2n/\log\log n)$ be the error parameter. The hash family $H^{\mathrm{md}}_{\beta,\delta}$ in Construction3 is locally explicit.
\end{authorclaim}
\begin{proof}
Recall that the hash function $H^{\mathrm{md}}$ in Construction3 is obtained by the 1-detector $L^{\mathrm{hg}}$ in Construction2 and the reduction in Construction1. The output of the hash $H^{\mathrm{md}}_{\beta,\delta}$ consists of two parts: the output of the 1-detector $L^{\mathrm{hg}}$, and several randomly sampled bits from the input. Since $L^{\mathrm{hg}}$ is locally explicit by Theorem3.12, it is sufficient to show that the sampling procedure is locally explicit. More precisely, given a seed $v$ of the sampling procedure and an index $i$, we need to compute efficiently all the indices $j$ such that $w_j=i$ for the vector $(w_0,\ldots,w_{b-1})$ corresponding to the seed $v$. This can be done using the algorithm $A^{\mathrm{samp}}$ from Lemma3.1.
\end{proof}
\begin{authorclaim}{Theorem3.14}
Let $\ell=\ell(n)=\operatorname{polylog}(n)$. The hash family $H^{\mathrm{expw}}_\ell$ in Construction4 is locally explicit.
\end{authorclaim}
\begin{proof}
Let $S_n=\{w_0,\ldots,w_{t-1}\}$ with $t=\widetilde O(n^2)$ is a $0.1$-biased set from Theorem2.4 and $G_t$ is the expander in Theorem2.10. The hash function in Construction4 is defined as
\[h_v(x)\triangleq\langle w_{v_0},x\rangle\Vert\langle w_{v_1},x\rangle\Vert\cdots\Vert\langle w_{v_{\ell-1}},x\rangle,\]
where $v=(v_0,\ldots,v_{\ell-1})\in[t]^\ell$ is a random walk on $G_t$. Given the seed $v$ (i.e., the seed of a random walk), we can produce $v_0,v_1,\ldots,v_{\ell-1}$ in $\operatorname{polylog}(n,\ell)=\operatorname{polylog}(n)$ time by Theorem2.10. Then using the algorithm $A$ in Theorem2.4 we can check for each $j$ whether the $i$th bit of $w_{v_j}$ is 1, which indicates whether the $j$th output bit is influenced by the $i$th input bit. Enumerating all $j\in[\ell]$ gives a required $\operatorname{polylog}(n)$ time algorithm listing all the outputs influenced by a particular input bit.
\end{proof}
Using Theorem3.13, Theorem3.14, and Proposition3.11, we immediately obtain the explicitness of Construction5.
\begin{authorclaim}{Corollary3.15}
Let $\beta\in(0,2]$ be a parameter and $\ell=\ell(n)\le O(\log^2n/\log\log n)$ be a non-decreasing function. The hash family $H^{\mathrm{final}}_{\beta,\ell}$ in Construction5 is locally explicit.
\end{authorclaim}
\subsection*{3.5 Uniformity of our construction}
We have shown in Theorem3.9 that our hash function can be computed by extremely sparse $\mathsf{CC}^0[2]$ circuit or $B_2$ circuits with small size and depth simultaneously. It remains to clarify the uniformity of our hash construction, i.e., the complexity we need to construct the circuit given the input length and the seed.
The definitions of uniformity vary with respect to the complexity measures, which can be chosen according to the applications. Typical choices include space complexity (e.g., LOGSPACE-uniform, see Section6.2.1 of [3]), time complexity (e.g., polynomial-time uniform), parallel time (see, e.g., Section3 of [29]), and descriptive complexity (see, e.g., [20]). If we choose the time complexity, for instance, we can define P-uniformity as the existence of the polynomial-time algorithm that prints the circuit given $1^n$ and the seed $v$. Indeed, it is easy to check that the $2n+o(n)$ size $B_2$ and $\mathsf{CC}^0[2]$ circuits for Construction5 in Theorem3.9 is P-uniform.
Since our hash function can be evaluated in low-depth circuit models such as $\mathsf{NC}^1$ and $\mathsf{CC}^0[2]$, the parallel time to generate the circuit also seems to be crucial for its further applications. Therefore in this section we will discuss the POLYLOGTIME-uniformity of our hash functions, defined as follows.
\begin{authorclaim}{Definition3.16}
Let $C$ be a circuit class. A family of $\operatorname{polylog}(n)$-samplable distributions over $m$-output functions (e.g., hash functions or 1-detectors) $H=\{H_n\}_{n\ge1}$ is said to be computable with POLYLOGTIME-uniform $S(n)$-size (measured in size, depth, wire complexity, etc) $C$-circuits if each of the functions $h_v\in\operatorname{supp}(H_n)$ corresponding to the seed $v$ is computable by an $S(n)$-size $C$-circuit, supplemented with the following $\operatorname{polylog}(S(n))$-time algorithms:
\begin{itemize}
\item $\mathrm{SIZE}(n,v)$ returns the size of the $C$-circuit $C_v$ computing $h_v$;
\item $\mathrm{TYPE}(n,v,i)$ returns the type of the $i$th node in the circuit $C_v$, which indicates (1) whether it is a gate, an input variable or a constant; (2) whether it is an output node; and (3) the type (if it is a gate or a constant) or index (if it is an input variable and/or output node) of it;
\item $\mathrm{EDGE}(n,v,i,j)$ returns the $j$th input of the $i$th node if the $i$th node is a gate in $C_v$.
\item $\mathrm{OUT}(n,v,i)$ returns the number of the $i$th output node.
\end{itemize}
It is guaranteed that the nodes are numbered in topological orders, and in particular, the input variables are numbered from 1 to $n$.
\end{authorclaim}
We claim that all our results of circuit complexity for the constructions in the previous parts of this section are in fact POLYLOGTIME-uniform. It is straightforward to check the claim, so we only sketch the proof and left the details to the readers.
\begin{authorclaim}{Theorem3.17 (Uniformity of Theorem3.9)}
Let $\beta\in(0,2]$ be a parameter and $\ell=\ell(n)\le O(\log^2n/\log\log n)$ be a non-decreasing function. The hash function $H^{\mathrm{final}}_{\beta,\ell}$ in Construction5 can be computed either by POLYLOGTIME-uniform depth-$(1+\lceil2/\beta\rceil)$ $\mathsf{CC}^0[2]$ circuits of wire complexity $2n+O(n/\log\log n)$, or by POLYLOGTIME-uniform $B_2$ circuits of size $2n+O(n/\log\log n)$ and depth $\log n+O(\beta^{-1})$ simultaneously.
\end{authorclaim}
\paragraph{Sketch.}
Since the POLYLOGTIME-uniformity is close under composition, we only need to verify that Construction3 and Construction4 are POLYLOGTIME-uniform. The uniformity of Construction3 directly follows from the algorithm $B$ in Theorem2.12, the algorithm $A$ in Theorem2.8, and the algorithm $A$ in Lemma3.1. The uniformity of Construction4 follows from the algorithm $A$ in Theorem2.4 and the algorithm $A$ in Theorem2.10. Note that the circuit for Construction4 should be slightly modified to make $\mathrm{EDGE}(n,\cdot,\cdot)$ $\operatorname{polylog}(n)$-time computable. Take the $\mathsf{CC}^0[2]$ case for example: each of the output gates has fan-in exactly $n$ (instead of $|w_{v_i}|$ for the $i$th output bit, see Construction4 for the notation); the $j$th input wire of the $i$th output gate is connected to the $j$th input or a constant 0 according to the $j$th bit of $w_{v_i}$. \hfill$\square$
